\chapter{Introduction}\label{ch01:intro}

\section{What this monograph is about}\label{ch01:about}

This monograph develops, at full technical resolution, a single candidate object: CAS-Ohm 1.0, a profile-specific experimental response and calibration model for steered neural-network activations. The candidate proposes that, within a narrowly declared experimental regime, a declared intervention coordinate --- the steering alpha applied along a fixed direction vector at a fixed test point in a fixed model --- and an observed activation coordinate --- the directional projection of the resulting activation --- may stand in a local, approximately proportional relation, and that the proportionality parameter, once fitted on a registered calibration split and carried with uncertainty, may predict held-out response. The candidate prediction is
\[
\hat{y} = \frac{a}{r_{\mathrm{eff}}},
\]
where $a$ is the alpha dose, $\hat{y}$ is the predicted activation response, and $r_{\mathrm{eff}}$ is the effective resistance in model-specific units of alpha per activation.

Stated so briefly, the proposal sounds like a modest exercise in curve fitting, and in one sense it is exactly that. The reason it warrants a monograph is not the arithmetic but the discipline surrounding the arithmetic. The historical source material from which CAS-Ohm 1.0 descends --- a quantification guide, an operationalization addendum, and an appendix to a printed-circuit-board-styled engineering text --- reached repeatedly for the vocabulary of Ohm's law: voltage, current, resistance, series and parallel composition, gain cascades, hidden resistance inferred from residuals. That vocabulary is seductive precisely because it comes bundled with an enormous inherited apparatus of intuitions, identities, and network theorems, almost none of which have been earned in the setting where the sources deploy them. A preserved source audit, reproduced in full later in this monograph, registers eight distinct conflicts or under-specifications across those three sources: mixed measurement scales, incompatible full-scale calibration methods, a missing dimensional bridge between score-valued and activation-valued quantities, two mathematically distinct definitions of gain, a mnemonic classification axis overloaded with two independent meanings, network composition rules asserted without validation, residual quantities pressed into causal service they cannot bear, and quantization band boundaries left without a precision rule.

CAS-Ohm 1.0 is the constructive response to that audit. It is a typed proposal in which every quantity carries its kind, unit, measurement provenance, calibration binding, and uncertainty; in which the historically observed failure modes are either unrepresentable in the type system or fail closed at validation; in which the Ohm-shaped prediction is admissible only under explicit anti-circularity and split-discipline conditions; and in which every empirical question the analogy raises is registered as a falsifiable hypothesis with a preregistered closure protocol rather than asserted as a fact. The candidate's own status vocabulary is blunt: at the time of writing, every CAS-Ohm claim is \textsc{open / ring 4 / not licensed / held}. Nothing in this monograph reports a validated empirical result about any live model, and nothing in it should be read as one.

This first chapter fixes the frame for everything that follows. Section~\ref{ch01:borrow} states precisely what CAS-Ohm borrows from Ohm's law --- the \emph{shape} of the relation and nothing else --- and Section~\ref{ch01:refuse} states, with equal precision, the four things it refuses. Section~\ref{ch01:terminology} sets out the terminology discipline that separates the candidate's quantities from physical electrical quantities. Section~\ref{ch01:dimensional} explains why dimensional consistency, though enforced everywhere in the candidate, is a necessary condition for meaningfulness and not a sufficient condition for empirical validity. Section~\ref{ch01:contributions} states the three contributions of the work. Section~\ref{ch01:charter} records the charter status of this document and the authority boundary that governs every claim in it. Section~\ref{ch01:guide} is a reading guide to the monograph.

\section{The borrowed shape}\label{ch01:borrow}

Ohm's law, in its elementary physical form, relates the potential difference across a conductor, the current through it, and a material- and geometry-dependent parameter called resistance. Three structural features of that relation --- and only those three --- are borrowed by CAS-Ohm.

\begin{definition}[The borrowed shape]\label{ch01:def:shape}
The \emph{borrowed shape} of Ohm's law consists of exactly three components:
\begin{enumerate}
\item a \emph{dose coordinate}: a scalar quantity that the experimenter declares and controls, swept over a registered grid;
\item a \emph{response coordinate}: a scalar quantity that the experiment observes, measured by a declared method under a declared calibration profile; and
\item a \emph{fitted proportionality parameter}: a scalar relating the two, estimated from calibration data by a declared fit method, carried with uncertainty, and evaluated on data disjoint from the data used to estimate it.
\end{enumerate}
\end{definition}

In the CAS-Ohm instantiation, the dose coordinate is the steering alpha $a$: the scalar coefficient by which a fixed, unit-norm, sign-declared direction vector is scaled before being added to (or otherwise composed with) the activation at a declared layer, token-position set, and hook type of a declared model. The response coordinate is the activation response $y$: the directional projection of the observed activation onto a declared effect direction at a declared test point. The fitted proportionality parameter is the effective resistance $r_{\mathrm{eff}}$, whose units are alpha per activation, so that the candidate prediction $\hat{y} = a / r_{\mathrm{eff}}$ is dimensionally coherent by construction.

The borrowed shape is deliberately austere. It does not include the physics of conduction, the microscopic theory that explains why some materials are ohmic over wide ranges, the thermodynamic account of dissipation, Kirchhoff's laws, the network theorems that follow from them, or the metrological infrastructure --- standard resistors, Josephson voltage standards, quantum Hall resistance standards --- that gives physical electrical measurement its extraordinary reproducibility. Every one of those omissions is intentional, and several of them are enforced by the type system rather than left to authorial restraint.

The intellectual tradition the borrowed shape actually belongs to is not circuit physics but dose-response modeling: the empirical study, ubiquitous in pharmacology, toxicology, and agronomy, of how a controlled stimulus maps to an observed response, with characteristic regions --- a sub-threshold region where response is indistinguishable from baseline, an approximately linear middle region, and a saturation region where further dose produces diminishing response \cite{ritz}. CAS-Ohm's dose-response region classification (which includes cutoff, active, saturation, breakdown, non-ohmic, and unstable regions, together with an unassessed default) descends from that tradition, not from electrical engineering. The Ohm vocabulary survives only as the name of the candidate and as the mnemonic for the fitted-slope-through-the-origin form of the active-region model; the analytical machinery is dose-response analysis under metrological discipline.

A second point deserves emphasis before any formalism appears: borrowing a shape licenses nothing. The relation $\hat{y} = a / r_{\mathrm{eff}}$ is a hypothesis schema, not a law. Whether any (model, layer, direction, test point, corpus, alpha grid) configuration actually exhibits a local linear response region; whether a fitted $r_{\mathrm{eff}}$ predicts held-out response better than trivial baselines; whether independent manipulation of proposed constraint mechanisms moves response with the predicted inverse slope --- these are the registered hypotheses of the falsification program, every one of which may fail. The monograph's stance is that the analogy is valuable exactly insofar as it is testable and rejectable, and worthless the moment it is treated as inherited truth.

\begin{example}[Schematic reading of the borrowed shape]\label{ch01:ex:shape}
To fix intuitions without asserting any result, consider a purely schematic run. An experimenter declares a calibration profile: an exact model identity (weights, configuration, and tokenizer digests), a layer, a token-position set, a unit-norm direction with declared sign and provenance, an alpha grid, a full-scale method, fit and uncertainty methods, and controls. Sweeping $a$ over the grid on the calibration split, the experimenter records typed observations of $y$. Suppose --- schematically --- the region classifier finds a contiguous sub-grid classified \textsc{active}. A fit on the calibration split produces $r_{\mathrm{eff}}$ with a standard uncertainty and interval, in role \textsc{calibrated held-out} bound to the split reference. Only now is $\hat{y} = a/r_{\mathrm{eff}}$ computable on the disjoint held-out split, and only its held-out performance --- against zero-intervention, random-direction, wrong-layer, and norm-matched controls --- would bear on the local-response hypothesis. Every step of this schematic is specification; whether any real configuration completes it with favorable numbers is precisely what no part of this monograph claims.
\end{example}

\section{The four refusals}\label{ch01:refuse}

The candidate is defined as much by what it refuses as by what it proposes. Four refusals are load-bearing; each is stated here as an explicit claim boundary and enforced later by the ontology, the audit dispositions, or the capability gates.

\begin{boundary}[No physical law]\label{ch01:bd:physical}
CAS-Ohm 1.0 asserts no physical law. No CAS-Ohm quantity is a physical voltage, current, resistance, charge, flux, power, or any other SI electrical quantity. No CAS-Ohm record is physical telemetry. The unit ontology contains no volt, ampere, or ohm; those labels are forbidden aliases for model quantities, and a hypothetical physical profile would require separate explicit typing, separate instrumentation, and separate validation that do not exist in the current candidate.
\end{boundary}

The first refusal is not a stylistic disclaimer but a structural fact about the type system. Chapter~\ref{ch03:ontology} presents a unit ontology whose members are \texttt{ALPHA}, \texttt{ACTIVATION}, \texttt{ALPHA\_PER\_ACTIVATION}, \texttt{DIMENSIONLESS}, and \texttt{FRACTION}, together with the \texttt{Q10} and \texttt{Q100} score scales. There is no representable way, within the candidate, to record a value in volts. The generic CAS PCBA vocabulary elsewhere in the wider project contains a member \texttt{Unit.VOLT} for entirely different purposes; the candidate's position, developed in Chapter~\ref{ch03:ontology}, is that the existence of that vocabulary member licenses nothing --- neither physical CAS telemetry nor a CAS-Ohm unit.

\begin{boundary}[No electrical implementation]\label{ch01:bd:electrical}
CAS-Ohm 1.0 describes no electrical circuit, device, board, or component, and no electrical behavior of any computing substrate. The series, parallel, and gain-composition operations defined in the candidate are structural research calculations over typed records, capability-gated and unauthenticated; they are not claims that any network of model components behaves like a resistor network.
\end{boundary}

The second refusal addresses the most tempting inheritance of all: network theorems. Series addition of resistances, reciprocal aggregation of parallel resistances, and multiplication of cascaded gains are, in circuit theory, consequences of conservation laws that hold with extraordinary generality in physical circuits. Nothing analogous has been established for steered activations. The source audit found the network rules presented in the Guide as working machinery, applied in Appendix~I with partial warnings, and correctly demoted in the addendum to preregistered predictions. The candidate adopts the strictest reading: composition operations exist only as gated structural calculations whose empirical content is exactly the registered composition and redundancy hypotheses, and no more.

\begin{boundary}[No attention measurement]\label{ch01:bd:attention}
CAS-Ohm 1.0 does not measure attention. No CAS-Ohm quantity is a measurement of a model's attention weights, attention allocation, or any cognitive construct named ``attention,'' and no CAS-Ohm record supports the product wording ``measures attention.''
\end{boundary}

The third refusal severs a specific historical thread. The source Guide computed a quantity it called raw attention as a ratio of score-valued quantities, inviting the reading that the framework measures how much attention a system pays to something. The candidate's response coordinate is nothing of the sort: it is a directional projection of an activation vector --- a geometric quantity relative to a declared direction, at a declared layer and token position, under a declared calibration profile. Whether such projections relate in any lawful way to attention mechanisms in the architectural sense, or to attention in any psychological sense, is outside the candidate's claim surface entirely.

\begin{boundary}[No causal law]\label{ch01:bd:causal}
CAS-Ohm 1.0 asserts no causal law. A fitted response relation is an observational-and-interventional summary within one calibration profile; it does not identify mechanisms, does not license counterfactuals beyond its registered regime, and does not support causal attributions from residual signs. Causal interpretation requires an explicit causal model and identifiable interventions \cite{pearl}, neither of which the candidate supplies.
\end{boundary}

The fourth refusal has teeth in a specific place: residual diagnostics. The sources repeatedly used the sign of a prediction residual to suggest hidden opposition, extra amplification, or hidden resistance --- in one case treating a positive residual on a safety-relevant configuration as evidence of a bypass. The candidate's residual discipline, developed later in the monograph, permits residuals to rank follow-up probes and nothing else: a residual never proves a hidden policy, a bypass, an opposition, a cause, or a safety failure, and at zero or zero-compatible predicted response the diagnostic ratios are censored rather than reported. The registered residual-localization hypothesis states the exact condition under which any stronger language could ever be earned --- blinded implanted faults localized above shuffled and matched controls --- and commits to withdrawing causal language if that test fails.

These four refusals interlock. A reader who grants the candidate its dose-response framing but silently retains an electrical mental model will over-trust the network operations; a reader who accepts the projection-based response coordinate but hears ``attention'' will over-read the semantics; a reader who accepts the fitted slope but forgets the causal boundary will convert residuals into accusations. The monograph therefore repeats the refusals wherever the corresponding machinery appears, at the cost of some redundancy, because the failure mode being defended against is precisely the quiet reactivation of inherited intuitions.

\section{Terminology discipline}\label{ch01:terminology}

The historical sources use $V$, $I$, and $R$ throughout, by analogy. This monograph and the candidate do not. The terminological rule is simple and total.

\begin{ruleenv}[Terminology discipline]\label{ch01:rule:terminology}
In every CAS-Ohm context, the dose coordinate is named the \emph{alpha dose} and written $a$; the response coordinate is named the \emph{activation response} and written $y$; the fitted proportionality parameter is named the \emph{effective resistance} and written $r_{\mathrm{eff}}$, with units of alpha per activation. The words \emph{volt}, \emph{ampere}, and \emph{ohm}, the symbols $V$, $I$, and $R$ in their electrical senses, and every derived electrical unit are reserved for physical profiles, which do not exist in the current ontology. Source quotations retain source notation, always inside an explicit attribution to the source.
\end{ruleenv}

The rule does more work than it may appear to. First, it prevents dimensional camouflage. When a score-valued quantity on a 0--10 scale is called $R$ and divided into a quantity called $V$, the notation itself suggests that the quotient is a current-like quantity, and the sources fell into exactly this trap: the audit's third conflict records that the predicted-response formula retained the form $V/R_{\mathrm{total}}$ even after $V$ had been reinterpreted as a steering alpha and $R$ remained score-derived, leaving the prediction dimensionally unmoored. Under the discipline, the same error is impossible to write down without noticing: a constraint score in \texttt{Q10} or \texttt{Q100} units simply is not an $r_{\mathrm{eff}}$ in \texttt{ALPHA\_PER\_ACTIVATION} units, and the type system rejects the division at the unit-validation layer.

Second, the rule keeps the analogy honest about its direction of travel. Physical Ohm's law earned its network theorems from conservation laws; the candidate's composition rules are hypotheses. Writing $r_{\mathrm{eff}}$ rather than $R$ is a persistent, low-cost reminder that the parameter is an effective, profile-relative, fitted quantity --- closer in spirit to an effective dose-response slope than to a material property.

Third, the rule fixes the reading of the one place where electrical vocabulary legitimately appears: the forbidden-alias list. The unit ontology names \textsc{volt}, \textsc{ampere}, and \textsc{ohm} precisely in order to forbid them. This is not decoration. An ontology that merely omitted the electrical units would be vulnerable to a later ``helpful'' extension adding them; an ontology that names them as forbidden aliases makes the refusal itself part of the contract, testable by the golden rejection fixtures that accompany the implementation.

A related discipline governs the word \emph{resistance} itself. The candidate uses \emph{effective resistance} only for the fitted \texttt{ALPHA\_PER\_ACTIVATION} parameter with a declared role, and \emph{constraint score} for the score-valued quantities the sources sometimes called resistance. The two are distinct quantity kinds with disjoint allowed units, and the audit's binding disposition --- constraint scores never masquerade as effective resistance --- is enforced at validation. Where the monograph must discuss the sources' conflated usage, it does so in quotation and with the conflict explicitly cited.

\section{Dimensional consistency is necessary, not sufficient}\label{ch01:dimensional}

Everything in the candidate is dimensionally disciplined: every quantity kind has a closed set of allowed units, every uncertainty carries the unit of its measurement, and the prediction $\hat{y} = a/r_{\mathrm{eff}}$ type-checks only because $r_{\mathrm{eff}}$ carries alpha-per-activation units. It is therefore worth stating plainly what this discipline does and does not buy.

Dimensional analysis, in the tradition formalized by Buckingham \cite{buckingham}, establishes that physically meaningful relations must be expressible in terms of dimensionless combinations of the governing quantities, and that a dimensionally inconsistent equation cannot be a correct description of anything. The contrapositive is the useful direction here: dimensional consistency is a \emph{necessary} condition for a relation to be a candidate description at all. An equation that adds a score to a projection, or divides an alpha by a 0--100 constraint score and calls the result an activation, is not wrong about the world; it is not yet about the world. The audit's mixed-scale conflict --- the worked example in which a saturation threshold of $8.5$ on a 0--10 scale was compared against region fractions on a 0--1 scale, with the corrected component value being $0.85$ --- is exactly such a pre-empirical failure, and the candidate's separation of \texttt{Q10}, \texttt{Q100}, and \texttt{FRACTION} units exists to make that class of failure unrepresentable.

But necessity is not sufficiency, and the distinction is the hinge of the entire claim discipline. A dimensionally impeccable relation can still be empirically false, vacuous, or circular. Three specific insufficiencies matter for CAS-Ohm.

First, \emph{fit does not imply prediction}. Given any paired sample $(a, y)$ with $y \neq 0$, the same-sample identity $r = a/y$ followed by $\hat{y} = a/r$ reproduces $y$ exactly, with perfect dimensional hygiene and zero empirical content. The candidate rejects this circularity by construction: an effective resistance is admissible for prediction only in the role \textsc{calibrated held-out} (or a separately validated network-equivalent role), which requires a registered calibration split disjoint from the evaluation split, a declared fit method, and carried uncertainty. The apparent-same-sample role exists in the ontology precisely so that the circular quantity can be recorded, labeled as what it is, and barred from prediction.

Second, \emph{units do not certify measurement}. A value in \texttt{ACTIVATION} units is not thereby a trustworthy observation; it might be a design declaration, a model-derived computation, or a score. The measurement-kind axis of the ontology --- design value, instrumented, behavioral proxy, model derived, human score, LLM score --- is orthogonal to the unit axis, and the non-fungibility rules (a derived value never becomes instrumented; a score never becomes telemetry) prevent dimensional correctness from laundering provenance. This is elementary metrological discipline \cite{vim,gum}, adopted without any claim that CAS-Ohm's instruments are themselves validated.

Third, \emph{dimensional coherence does not select the model}. Within the same unit system one can write a proportional law, a saturating law, a threshold law, or an arbitrary monotone curve; dimensional analysis is silent among them. Which, if any, describes a given configuration is an empirical question for the falsification program, to be answered under calibration discipline that represents parameter uncertainty and model discrepancy rather than anointing a best fit as truth \cite{kennedyohagan}, and under the separation of software verification from operational validation \cite{sargent}.

The monograph therefore treats dimensional consistency as a floor: a precondition enforced everywhere, celebrated nowhere. Any prose in later chapters that appears to move from ``the units work out'' toward ``the relation holds'' should be read against this section and rejected; no such move is licensed.

\section{The three contributions}\label{ch01:contributions}

The monograph makes three contributions. None is an empirical finding about any live model; all are artifacts of specification, preservation, and program design.

\subsection{An ontology in which the observed failure modes are unrepresentable}\label{ch01:contrib:ontology}

The first contribution is the quantity and evidence ontology of Chapter~\ref{ch03:ontology}: a typed system in which every numeric record carries its quantity kind (one of sixteen), its unit (from a closed unit ontology with forbidden electrical aliases), its measurement kind (one of six, non-fungible), its method identity, its calibration-profile digest, its estimate-or-observation status, its uncertainty and confidence interval, its sample size, its raw-artifact and control references, its derivation parents, and --- where applicable --- its missing-or-censored reason. No bare binary float is admitted anywhere in the evidence layer.

The design criterion for the ontology was neither generality nor elegance but a specific negative property: each failure mode surfaced by the source audit should be either unrepresentable or fail-closed. The mixed-scale comparison cannot be written because \texttt{Q10}, \texttt{Q100}, and \texttt{FRACTION} are distinct units with disjoint validation ranges and no silent conversion. The dimensionally unmoored prediction cannot be written because effective resistance is the only quantity licensed to bridge alpha and activation, and it demands its own units, role, and split binding. The score-as-telemetry slide cannot occur because measurement kinds are non-fungible and instrumented records demand instrument identity, calibration binding, raw evidence, and sample information. The residual-as-cause leap is blocked because residual emission requires an admitted active-region prediction with propagated uncertainty and is censored at zero-compatible response. The ontology is, in this sense, the audit's negative image.

\subsection{A preserved audit at evidence-record granularity}\label{ch01:contrib:audit}

The second contribution is the preserved source audit itself. Eight conflicts or under-specifications across the three sources are registered in a machine-readable evidence record, each with disclosed locator systems (rendered PDF pages for the Guide; OOXML paragraph and table-row indices for the addendum, for which no page number is asserted because both available page-rendering routes failed; UTF-8 line numbers for Appendix~I; PDF and book pages for the PCBA gain definitions), the conflicting facts side by side, and a binding disposition for the candidate.

Preservation, rather than harmonization, is the point. It would have been easy --- and it would have been malpractice --- to smooth the sources into a single retrospectively coherent story. The preserved audit makes that impossible: downstream code and prose cannot claim one coherent inherited scale, one full-scale method, or one gain law where the sources demonstrably provide several, and every disposition in the candidate is traceable to the specific conflict it resolves or quarantines. The audit is thus an engineering result in its own right: it converts historical incoherence from a latent hazard into an explicit, cited, machine-checkable boundary condition.

\subsection{A fail-closed candidate with a falsification program}\label{ch01:contrib:candidate}

The third contribution is the candidate itself: a specified, implemented, and test-covered system --- calibration-profile contract; independent dose-response-region and semantic-outcome classification axes; exact decimal quantization with a boundary-preservation guarantee; capability-gated structural network operations; residual censoring; and a strict split between structural research-packet completeness and authenticated Evidence-Packet eligibility, the latter constant false in this candidate --- together with a falsification program of five registered hypotheses (local response, effective-resistance prediction, gain composition, redundancy, residual localization) and a preregistered closure protocol fixing controls, falsifiers, negative-outcome retention, and permitted claim language before any run.

The phrase \emph{fail-closed} carries specific weight. Wherever the candidate encounters a record that is score-only, zero-instrumentation, missing-provenance, mixed-scale, unknown-profile, or tampered, it rejects rather than degrades; wherever a capability gate's conditions are unmet, the operation is unavailable rather than approximate; wherever authentication conditions are unmet --- which, in this candidate, is always --- eligibility is false rather than provisional. The bounded software evidence that exists (a passing test suite exercising rejection paths, boundary cases, and gate behavior at a pinned receipt) establishes tested software behavior only, and the monograph will say so every time it is mentioned.

\section{Charter status and the authority boundary}\label{ch01:charter}

\subsection{Provisional technical report, not an eighth charter paper}\label{ch01:charter:status}

The Project Coherence charter, which governs the wider research program this work sits inside, lists seven planned papers. No standalone CAS-Ohm paper is among them. Until and unless a charter amendment exists, this monograph and the research paper it extends are \emph{provisional technical reports}: internal working artifacts, drafted to fix the candidate's specification and claim discipline at full resolution, and explicitly not representable as an eighth charter paper, as a charter-approved publication, or as any other charter-level deliverable.

This is not modesty theater; it is a governance fact with operational consequences. Charter papers pass through the program's admission machinery --- internal cross-review, independent external review, claim-ledger reconciliation --- before any claim in them changes status. This monograph has passed through none of that: drafting it promotes no claim, and both internal cross-review and independent external review remain open. A reader who cites this monograph should cite it as what it is: a draft technical report describing a candidate whose every substantive claim is registered as open.

\subsection{The authority boundary}\label{ch01:charter:authority}

\begin{boundary}[Authority boundary]\label{ch01:bd:authority}
This monograph establishes no live execution, no scientific conclusion, no physical measurement, no causal effect, no safety or clinical conclusion, no product capability, no license, no canon admission, and no disclosure. All CAS-Ohm claims referenced herein hold status \textsc{open / ring 4 / not licensed / held}. Authenticated Evidence-Packet eligibility is constant false in this candidate. Structural consistency, canonical digests, replayed attestations, and passing test suites are integrity and replay evidence only; none is reviewer, legal, product, or scientific authority.
\end{boundary}

The boundary deserves a paragraph of unpacking, because its final sentence names the subtlest hazard. The candidate is unusually rich in integrity machinery: canonical byte encodings, SHA-256 digests over every typed record, optional Ed25519 fixture-key attestations that are verified and replayed, exhaustive lineage. It is natural --- and wrong --- to let that machinery accrete into an impression of authority. A digest proves that bytes have not changed; it does not prove the bytes were ever true. A replayed attestation proves a fixture key once signed a payload; it does not conjure an independent reviewer. Provenance structure is not truth or authority \cite{prov}. The candidate encodes this as an architectural split --- completeness and eligibility are independent axes, and no accumulation of structural evidence can move a packet from complete to authoritative --- and the monograph adopts the same split in its prose: descriptive thoroughness in the chapters that follow is never to be read as an authority claim.

A final component of the boundary concerns forbidden product wording. The claim ledger records, as a \textsc{not licensed} product claim, any blanket representation of design scores, behavioral proxies, or model-derived values as physical telemetry; and the specification forbids, in any product context, the phrases ``physical telemetry,'' ``measures attention,'' ``electrical law,'' ``causal by construction,'' ``cross-model comparable,'' ``proves hidden resistance,'' and any safety or verdict claim. The monograph mentions these phrases only to forbid them.

\section{Relation to the sources and to the wider program}\label{ch01:relation}

\subsection{The three sources}\label{ch01:relation:sources}

Three documents constitute the audited lineage. The \emph{Quantification Guide} is the earliest and most ambitious: it introduces the electrical vocabulary, defines score-valued drive and constraint quantities on a 0--10 scale (constraint on $0.1$--$10.0$), computes a raw attention quantity as their ratio, presents series, parallel, and gain-cascade rules as working machinery, reads residual signs diagnostically, and displays a four-band quantization table with two-decimal band edges. The \emph{operationalization addendum} is the pivotal reinterpretation: it maps the dose coordinate to a steering alpha and the response coordinate to a directional activation projection --- the move that makes the proposal empirically testable at all --- while retaining, in places, score-derived resistance components inside the predictive formula, and it is also the source that correctly demotes the network rules to preregistered predictions and separates dose-response signatures from effect-axis outcomes. \emph{Appendix~I} of the PCBA-styled engineering text is the most self-critical source: it rescales constraint to 0--100, moves saturation and cutoff to 0--1 fractions, identifies the Guide's mixed-scale worked example as an error, defines a per-path rail-and-floor full-scale procedure and warns it does not transfer, frames the predictive relation as testable design arithmetic rather than physics, and warns that co-present guardrails are not parallel routes --- while nonetheless itself treating a positive residual on a safety board as a bypass.

The candidate's relation to these sources is neither adoption nor repudiation but audit-mediated inheritance. Where the sources agree and the content survives typing, the candidate adopts it. Where they conflict, the conflict is preserved as a first-class record and the candidate binds a disposition that is deliberately the most conservative reading --- distinct quantity types over silent conversion, separately named calibration methods over a preferred one, preregistered predictions over inherited theorems, censored residuals over diagnostic ones. Where a source under-specifies, the candidate specifies exactly and marks the specification as its own choice rather than as recovered intent. Chapter~\ref{ch03:ontology} and the audit chapter make each of these moves explicit; nothing is harmonized silently.

\subsection{Position inside the program}\label{ch01:relation:program}

Within the wider research program, CAS-Ohm 1.0 occupies a deliberately peripheral position. It inherits the program's canonical-encoding contract (the evidence bounds of Chapter~\ref{ch03:ontology} are an inheritance, not an invention), its claim-ledger discipline, its ring-based admission model, and its separation of structural consistency from authorization. It contributes back, so far, only working material: a typed candidate, a preserved audit, and registered hypotheses. It depends on no result from the program's other tracks, and no other track depends on it; a total failure of the CAS-Ohm falsification program would falsify the Ohm-shaped proposal and nothing else. This isolation is intentional. A candidate whose sources required an eight-conflict audit is exactly the kind of object that should be quarantined until its own program closes, and the monograph's authority boundary (Section~\ref{ch01:charter}) is the textual form of that quarantine.

\subsection{How to disagree with this monograph}\label{ch01:relation:disagree}

A specification-first document invites two distinct kinds of disagreement, and the claim discipline benefits from keeping them separate. The first kind targets the specification: a reader may hold that a binding disposition resolves a source conflict the wrong way, that a field of the observed-value record is redundant or missing, that a capability gate is too strict or too lax, or that a registered hypothesis is malformed. Such disagreements are welcome and actionable --- they are exactly what the open internal cross-review is for --- and because every disposition cites its conflict and every gate cites its rationale, a critic can attack the precise joint rather than the gestalt. The second kind targets the empirical question: a reader may expect the local-response hypothesis to fail, or the composition rules to be artifacts. Such expectations are not disagreements with this monograph, which asserts neither the hypotheses nor their negations; they are predictions about the falsification program's outcome, and the preregistered protocol is designed so that either outcome, once obtained under controls, will be recorded without erasure. The one form of engagement the monograph forecloses is treating its descriptive machinery as if it carried empirical or product authority --- in either direction.

\section{Status vocabulary and claim discipline}\label{ch01:status}

Because every substantive statement in this monograph is qualified by a status, the status vocabulary itself must be defined before it is used further. The vocabulary is inherited from the wider program's claim-ledger discipline and applied here without modification.

\subsection{The status line}\label{ch01:status:line}

Each CAS-Ohm claim carries a four-part status. \textsc{Open} states the research status: the claim is registered, unresolved, and awaiting evidence; nothing about it has been accepted or rejected. \textsc{Ring 4} states the admission ring: the outermost ring of the program's admission model, in which artifacts are working material --- specifications, candidates, drafts --- that have crossed none of the review gates separating working material from admitted results. \textsc{Not licensed} states the license status: no claim may be used to support a product statement, a marketing statement, a safety statement, or any external representation. \textsc{Held} states the disclosure status: the material is not released, and drafting or circulating it internally effects no disclosure.

The four parts are independent axes, not degrees of a single scale, and the candidate's implementation embeds all four constants into the canonical payload of every measurement record, so that the status travels with the bytes: any future consumer of a CAS-Ohm record will find \texttt{OPEN}, admission ring \texttt{4}, \texttt{NOT LICENSED}, and \texttt{HELD} inside the digested content itself, not in detachable metadata. A record cannot be quoted out of its status without altering its digest.

\subsection{What drafting does and does not do}\label{ch01:status:drafting}

A recurring hazard in long-form technical writing is that the act of careful exposition comes to feel like validation: a mechanism described in detail across many pages acquires, in the reader's mind, the solidity of a mechanism that works. The claim discipline addresses this directly. Drafting this monograph promotes no claim. Describing the anti-circularity rule does not establish that any non-circular fit succeeds; describing the capability gates does not establish that any capability's conditions have ever been met on live data; describing the falsification program does not establish that any hypothesis will survive it. The monograph's descriptive completeness is an input to future review, not a substitute for it.

Symmetrically, the discipline constrains the author as much as the reader. No sentence in this monograph may assert, of any live model, that a linear response region exists, that a fitted effective resistance predicted anything, that a network composition rule held, or that a residual localized a fault. Where the temptation to write such a sentence arises --- and in expository writing about a response model it arises constantly --- the licensed forms are the conditional (``if the local-response hypothesis survives its controls, then\ldots''), the registered (``the redundancy hypothesis states\ldots''), and the negative (``no such result exists at the time of writing''). The reader is invited to audit the monograph against this rule; any violation found is a defect of the text, not a leak of hidden evidence.

\subsection{Bounded software evidence}\label{ch01:status:software}

One category of evidence does exist and is cited where relevant: bounded software evidence. At the pinned receipt of the current gate, the candidate's test suite passes --- sixty-five test methods and fifty-eight parameterized subtests in the combined focused receipt --- exercising quantity and unit rejection, canonical-evidence limits, calibration identity, split separation, exact quantization boundaries, structural network constraints, residual censoring, uncertainty propagation, lineage, and fail-closed capability behavior, including synthetic known-fault boards (linear, cutoff, saturation, nonlinear, parallel, confounded, bypass) constructed to verify that the software detects what it claims to detect.

The scope of this evidence is exactly its name: it establishes tested software behavior, at a pinned revision, on synthetic inputs. It is verification, not validation \cite{sargent}. A suite that proves the classifier assigns band boundaries by exact decimal comparison proves nothing about whether any real configuration occupies a linear region; a synthetic bypass board that the residual machinery flags proves nothing about bypasses in any live model. The monograph cites the software evidence often, because fail-closed behavior is a genuine engineering property worth establishing; it converts that evidence into empirical support for the response model nowhere, because no such conversion is valid.

\section{Methodological stance: specification before experiment}\label{ch01:stance}

It is reasonable to ask why a candidate with no admitted empirical results warrants a specification of this length. The answer is the monograph's central methodological commitment: for a proposal whose historical sources demonstrably conflated scales, units, methods, and claim types, \emph{specification must precede experiment}, because an experiment run against an ambiguous specification produces evidence about nothing in particular.

Consider what an alpha-sweep experiment inherits from the specification layer. Which full-scale method normalizes the response --- a board-declared constant, a per-path rail-and-floor measurement, or a corpus percentile --- changes every normalized number downstream; the audit found all three in the sources, incompatible and unmarked. Whether a saturation threshold lives on a 0--10 scale or a 0--1 scale decides whether a worked example is in-region or seven times out of range; the audit found exactly that discrepancy. Whether gain means a Fisher--Rao displacement ratio or a local Jacobian projection derivative decides what a cascade experiment even measures; the audit found both, defined in different sources, with no bridge. An experiment run before these questions are pinned does not test the response model; it tests the experimenter's implicit choices, and its results cannot be compared, replicated, or even clearly stated.

The candidate therefore front-loads disambiguation. Every choice the sources left open is either fixed by a binding disposition (with the conflict preserved so the fix is visible as a fix), parameterized into the calibration profile (so that incompatible choices become distinct profiles with distinct digests rather than silent variants), or excluded by the type system. Only after that layer is stable does the falsification program define runs --- and by then, each registered hypothesis can state exactly which profile, grid, controls, and fit methods it binds, which is what makes its eventual outcome interpretable and its preregistration meaningful \cite{kennedyohagan}.

This stance has a cost, which the monograph accepts explicitly: it defers every interesting empirical question. The reader who wants to know whether steered activations are locally ohmic will not learn it here, and the author does not know either. What the reader will obtain is the precise, typed, falsifiable form of the question --- which, for a research lineage that previously could not state the question without dimensional error, is the necessary first result.

\section{Notation and conventions}\label{ch01:notation}

The following conventions hold throughout the monograph.

\emph{Quantities.} The alpha dose is $a$, in \texttt{ALPHA} units. The activation response is $y$, in \texttt{ACTIVATION} units; a predicted response is $\hat{y}$. The effective resistance is $r_{\mathrm{eff}}$, in \texttt{ALPHA\_PER\_ACTIVATION} units, strictly positive, and always accompanied --- in any admissible use --- by a role, a fit method, a split binding, and an uncertainty. Residuals are $\delta_y = y_{\mathrm{obs}} - y_{\mathrm{pred}}$. Score-valued quantities on the 0--10 and 0--100 scales are written with their units (\texttt{Q10}, \texttt{Q100}) rather than with suggestive symbols.

\emph{Enumerations.} Implementation enumerations are reproduced in \texttt{typewriter} face with their exact member spellings, including underscores, as in \texttt{CALIBRATED\_HELD\_OUT} or \texttt{CORPUS\_PERCENTILE\_PER\_TEST\_POINT}. Small-caps forms such as \textsc{instrumented} are prose references to the same members; where the two could diverge, the typewriter form is authoritative.

\emph{Exactness.} All boundary-sensitive arithmetic in the candidate is exact decimal arithmetic, never binary floating point; the monograph's worked numbers are to be read as exact decimals. Digests are SHA-256 over canonical bytes; equality of records means equality of digests.

\emph{Profiles.} Every quantity is relative to a calibration profile, identified by its canonical digest. Where the profile is clear from context it is suppressed from the notation; nothing transfers across profiles without a stated comparison rule, and no suppressed profile subscript should be read as cross-profile generality.

\emph{Citations.} Numbered citations refer to the external literature. Source documents of the audit (the Guide, the addendum, Appendix~I, and the PCBA text) are referred to by those names, with locators, and are deliberately not treated as citable scientific literature: they are audited inputs, and their conflicting content is evidence about the lineage, not authority about the world.

\section{Reading guide}\label{ch01:guide}

The monograph is organized so that a reader can always locate the epistemic status of what is being read: background, specification, preserved conflict, or registered hypothesis.

\textbf{Chapter~\ref{ch02:background}} is a tutorial survey of the five literatures the candidate stands on: dose-response modeling, with its response regions, linear ranges, and saturation phenomena \cite{ritz}; metrology, with its vocabulary of quantities, units, measurement kinds, uncertainty, and traceability \cite{vim,gum}; the statistical calibration of computer models, including parameter uncertainty, model discrepancy, and train-versus-held-out discipline \cite{kennedyohagan}; verification and validation as distinct activities \cite{sargent}; the limits of causal inference from observational fits \cite{pearl}; activation steering and representation engineering as the empirical setting, at primer depth \cite{zourepe}; and provenance modeling \cite{prov}. Everything in that chapter is background: none of it is a CAS-Ohm novelty, and none of it smuggles an empirical claim about CAS-Ohm.

\textbf{Chapter~\ref{ch03:ontology}} presents the quantity and evidence ontology in full: the observed-value record field by field with the rationale for each field; all sixteen quantity kinds; all six measurement kinds with examples and the non-fungibility rules; the unit ontology with its forbidden aliases; and the canonical evidence bounds --- exact-type tagging, cycle awareness, unsafe-Unicode rejection, and the declared depth, node, container, text, output, and integer limits, with host-exception normalization.

Subsequent chapters, outside the scope of this volume's opening movement, develop the preserved source audit conflict by conflict; the calibration-profile contract and the response model with its anti-circularity rule; the two independent classification axes and the exact quantization scheme with its boundary-preservation proposition; the capability-gated structural network operations; residual censoring; the completeness-versus-eligibility split; the five registered hypotheses and the preregistered closure protocol; the memristive extension boundary; and the candidate's limitations, stated without softening.

Three reading conventions apply throughout. First, environments carry status: a \emph{definition} or \emph{ruleenv} states candidate specification; a \emph{conflict} preserves a source disagreement without resolving the underlying question; a \emph{hypothesis} is registered and open; a \emph{boundary} bounds claims and is normative for every reading of the surrounding text. Second, enumerated implementation values are typeset in \texttt{typewriter} face and reproduce the implementation exactly; where the specification prose and the implementation enumerate slightly differently (as with the dose-response regions, where the implementation retains a \texttt{BREAKDOWN} member alongside the audit's five-region summary), both are recorded and the discrepancy is stated rather than smoothed. Third, the status line \textsc{open / ring 4 / not licensed / held} is not boilerplate to be skimmed; it is the single most important sentence in the monograph, and every other sentence is subordinate to it.

With the frame fixed, the monograph now turns to background: the established literatures whose disciplines the candidate borrows, and against whose standards its falsification program will eventually be judged.
